% ------------------------------------------------------------------------------
% Bab II - Landasan Teori
% ------------------------------------------------------------------------------
\chapter{LANDASAN TEORI}

% ------------------------------------------------------------------------------
% Tinjauan Pustaka
% ------------------------------------------------------------------------------
\section{Tinjauan Pustaka}
Bagian kajian pustaka diawali dengan ringkasan hasil kajian pustaka (paper atau jurnal) terkait research question atau problem statement yang diajukan pada Bab Pendahuluan. 

Contoh pustaka prosiding \cite{doyen2014explicit}, jurnal \cite{gunawan2015hydrostatic}\cite{panda} dan buku \cite{toro2013riemann}. Atau dapat juga menggunakan dua pustaka atau lebih dalam \cite{gunawan2015hydrostatic,toro2013riemann}.

% ------------------------------------------------------------------------------
% Landasan Teori
% ------------------------------------------------------------------------------
\section{Landasan Teori}
Teori pendukung dapat berasal dari textbook, paper/jurnal, tesis, dan TA lain dengan memperhatikan penggunaan teknik acuan sebagai bentuk kejujuran akademik.

\subsection{Sub Topik 1}
Menurut paper Kentang \cite{Kentang}, persamaan SWE adalah sebagai berikut:
\begin{equation}\label{Pers1}
    a=b+U^{n+1}_{i+1}
\end{equation}
Persamaan (\ref{Pers1}) adalah contoh persamaan.

\subsection{Sub Topik 2}
\begin{equation}\label{nama-rumus}
    \int_0^1 \frac{f(x)}{g(x)}\ {\rm dx}=\sin x
\end{equation}

\begin{equation}\label{nama-rumus1}
   \alpha \times \beta =\gamma^{3\alpha}
\end{equation}

\subsection{Sub Topik Dst}
Berikut adalah contoh gambar:
\begin{figure}[h!]
    \centering
    \includegraphics[scale=0.3]{Aset/Logo_Institusi.png}
    \caption{Contoh Logo Institusi}
    \label{fig:my_label}
\end{figure}
Lihat \textit{pada} Gambar \ref{fig:my_label}.